\documentclass[12pt]{article}
\usepackage[a4paper,margin=2.2cm]{geometry}
\usepackage{amsmath,amssymb,amsthm,mathtools}
\usepackage{hyperref}
\usepackage{cleveref}

% 中文（xelatex）
\usepackage{fontspec}
\usepackage{xeCJK}

\usepackage{fontspec}
\usepackage{xeCJK}
\setCJKmainfont{Noto Serif CJK SC}

\title{Orthogonal Projection}
\date{}
\begin{document}
\maketitle


\section{向量的正交分解与投影}
\begin{flushleft}
    \begin{itemize}
        向量 $\vec{u} = \begin{bmatrix}
                x \\
                y
            \end{bmatrix}, \vec{v} = \begin{bmatrix}
                x_{1} \\
                y_{1}
            \end{bmatrix}, \vec{u},\vec{v} \neq \vec{0}$

        向量本质是一个空间的移动指令(复合的), 这个空间是一个网格 由基底决定, 对于 u, v 而言是标准基(e1,e2)

        $\vec{u} = xe1 + ye2, \; \vec{v} = x_{1}e1 + y_{1}e2$

        也就是说 $x, y, x_{1}, y_{1} $ 都是基于e1, e2 的移动分量。\\
        现在我们设 $\vec{u} = X$, 这 X 是 无关基底的绝对地址。 \\

        在标准基网格环境, 通过  $xe1 + ye2 \to X $ 表达  \\
        那么我们设 在一组新的基底 $e1^{'}, e2^{'}$ 和新分量 $x_{new}, y_{new}$  表达\\
        存在下列代数组关系：
        \begin{align}
            xe1 + ye2 \to X \\
            x_{new}e1^{'} + y_{new}e2^{'} \to X
        \end{align}

        step1: $e1^{'}, e2^{'}$ 由 $\vec{v}$ 派生(方向归一化的正交分解), 基于 v 诞生的新基底。\\
        step2: $x_{new}e1^{'}, y_{new}e2^{'}$ 与 $xe1,ye2$ 的代数逻辑与几何投影关系。\\


        \vspace{0.5em}
        \item 1: $\vec{v}$ 方向归一化的正交分解($e1^{'},e2^{'}$) \\

              \vspace{0.3em}
              \textbf{1.1 $e1^{'}$} \\
              $e1^{'}$ 非常清楚和简洁 $\to$  $\vec{v}$ 的方向并且长度为1的单位 \\
              \begin{align}
                  e1^{'} = \frac{\vec{v}}{||\vec{v}||} \\
              \end{align}
              (3): 保留方向并归一化


              \vspace{0.3em}
              \textbf{1.2 $e2^{'}$} \\
              $e2^{'}$ 是 $e1^{'}$ 逆时针旋转 $90^{\circ}$  \\
              \begin{align}
                  e2^{'} = R_{90^{\circ}}\frac{\vec{v}}{||\vec{v}||}
              \end{align}
              (5): $R_{90^{\circ}} = \begin{bmatrix}
                      0 & -1 \\
                      1 & 0
                  \end{bmatrix} \to $逆时针旋转90

              \vspace{0.3em}

              逆时针旋转$90^{\circ}$ = $(e1^{'} \cdot e2^{'} = 0)$ = 正交性(这是代数和几何的统一, 不需要证明) \\
              到这里我们完成了 对$\vec{v}$ 基于归一化的正交分解, 获得的基于 v 的新世界的基底

              \vspace{0.5em}
        \item 2: $\vec{u} \;\text{与}\; \vec{v} $ 的代数逻辑与几何投影关系  \\
              1: 怎么基于 $e1^{'},e2^{'}$ 走到 X ? \\
              \begin{align}
                  \vec{u} = X = x_{new}e1^{'} + y_{new}e2^{'} \\
                  \vec{u} =  x_{new}e1^{'} + y_{new}e2^{'}
              \end{align}
              这里我们需要解决的问题是 算出$x_{new}, y_{new}$ 两个分量 \\
              已知正交性: \\
              \begin{align}
                  e1^{'} \cdot e2^{'} = 0 \\
                  e1^{'} \cdot e1^{'} = 1 \\
                  e2^{'} \cdot e2^{'} = 1 \\
              \end{align}
              $e1^{'}, e2^{'}$ 长度为1 \\
              存在代数转换: \\
              \begin{align}
                  \vec{u} =  x_{new}e1^{'} + y_{new}e2^{'}                             \\
                  \vec{u} \cdot e1^{'} =  (x_{new}e1^{'} + y_{new}e2^{'}) \cdot e1^{'} \\
                  \vec{u} \cdot e1^{'} =  x_{new}1 + y_{new}0                          \\
                  x_{new} = \vec{u} \cdot e1^{'}
              \end{align}

              \begin{align}
                  \vec{u} =  x_{new}e1^{'} + y_{new}e2^{'}                             \\
                  \vec{u} \cdot e2^{'} =  (x_{new}e1^{'} + y_{new}e2^{'}) \cdot e2^{'} \\
                  \vec{u} \cdot e2^{'} =  x_{new}0 + y_{new}1                          \\
                  y_{new} = \vec{u} \cdot e2^{'}
              \end{align}
              代入: \\
              \begin{align}
                  x_{new} = \vec{u} \cdot \frac{\vec{v}}{||\vec{v}||} \\
                  x_{new} = \frac{\vec{u} \cdot  \vec{v}}{||\vec{v}||}
              \end{align}
              到这里我们自然发现了 $x_{new} $ 它是 u 在  v 的标量投影 (这是完全基于换基视角"发现"的) \\
              向量投影(标量投影 * 单位向量): \\
              \begin{align}
                  x_{new}e1^{'} = (\frac{\vec{u} \cdot  \vec{v}}{||\vec{v}||})\frac{\vec{v}}{||\vec{v}||} \\
                  x_{new}e1^{'} = \frac{\vec{u} \cdot  \vec{v}}{(||\vec{v}||)^{2}}\vec{v}
              \end{align}
              也就是$x_{new}e1^{'}$ 就是 u 在 v 上的向量投影 = $proj_{\vec{v}^{\vec{u}}}$ \\

              \vspace{0.5em}
        \item 3 从内积公式代数出发
              内积可以整理称投影形式： \\
              $\vec{u} \cdot \vec{v} = |\vec{u}||\vec{v}| \cos\theta$ \\
              给定向量 u 与方向 v, 把 u 沿着 v 方向分解为两部分 $\to$  一部分 平行于 v, 一部分垂直于 v \\
              \begin{align}
                  \vec{u} = \vec{u}_{\parallel} + \vec{u}_{\perp}  \\
                  \vec{u}_{\parallel} = c\vec{v}                   \\
                  \vec{u}_{\perp} =  \vec{u} - \vec{u}_{\parallel} \\
                  <\vec{u}_{\perp}, \vec{v}> = 0                   \\
              \end{align}
              存在上述代数式关系 \\
              这里我们要弄清楚 $cv$ 是什么, 表示什么东西 \\
              c 是一个标量系数, 已知 v 是方向向量 那么 $cv$ 是对这个 v 做系数拉伸 \\
              我们知道在一组基底环境下(无关维度) 两个向量是满足平行四边形法则(构成一个平面) \\
              例如 u, v  会存在有关 v 系数的向量 $cv$ 会与 u 做垂线 构成直角三角形(平行四边形的一半) \\
              这个 $cv$ 也称为 $\vec{u}$ 在 $\vec{v}$ 方向上的正交投影向量 \\
              $proj_{v}(u) = cv = \vec{u}_{\parallel}$ \\
              那么算出 c 我们就可以得到 $proj_{v}(u)$ \\
              \begin{align}
                  <\vec{u}_{\perp}, \vec{v}> = 0                   \\
                  <\vec{u} - \vec{u}_{\parallel}, \vec{v}> = 0     \\
                  <\vec{u} - cv, \vec{v}> = 0                      \\
                  <\vec{u},\vec{v}> - c<\vec{v}, \vec{v}> = 0      \\
                  c = \frac{<\vec{u},\vec{v}>}{<\vec{v}, \vec{v}>} \\
              \end{align}
              $proj_{v}(u) = cv = \frac{<\vec{u},\vec{v}>}{<\vec{v}, \vec{v}>}\vec{v}$
              当 v 向量已经归一化($|v| = 1$) \\
              $proj_{v}(u) = <\vec{u},\vec{v}>\vec{v}$ \\

              这里有几个很有意思的点  \\
              1:  当 $|v| = 1$ 时  \\
              $proj_{v}(u) = <\vec{u},\vec{v}>\vec{v}$  \\
              也就是说 $ c = <\vec{u},\vec{v}>$, 那么点积到底是什么呢? 它把 u 常量化 和 v 是一种纯粹的倍数关系 \\
              2: $\vec{u}_{\parallel}, \vec{u}_{\perp}$ 和 $\vec{v}$ 的关系 \\
              它们使 $\vec{v}$ 的方向构造一个新网格(正交垂直), 这个网格和标准基的正交网格有区别吗?
              没有! 它只是看到这个世界的另一个视角 \\
              那么 $\vec{u} = c1\vec{u}_{\parallel} + c2\vec{u}_{\perp}$ 是完全成立的 \\
              $\vec{v}$ 被拆解成 一组分量 + 一组基 \\
              但$\vec{v}$ 在开始时 我们已经确了它只是一个方向, 也就是, 我们可以把一个环境转换成 $\to$ 任意视角"方向"的一个环境, 而且它们存在代数映射关系 ??

              因为我们是从内积公式出发的(双线性), 那么函数 f(x) 也是满足的: \\
              $proj_{g(x)}(f(x)) = \frac{\int f(x)g(x)}{(g(x))^{2}}g(x)$
    \end{itemize}
\end{flushleft}

\section{投影最小化距离}
\begin{flushleft}
    \begin{itemize}
        \item[] \textbf{定理}：在所有形如 $cv$ 的向量中，$proj_{v}(u)$ 是离 $u$ 最近的那个。即：
              $$ \min_{c \in \mathbb{R}} \|u - cv\| \text{ 在 } c = \frac{\langle u, v \rangle}{\langle v, v \rangle} \text{ 处取得} $$
              其中 $\frac{\langle u, v \rangle}{\langle v, v \rangle}$ 为投影系数。

              \vspace{0.5em}
        \item[] \textbf{证明}：
              对于任意实数 $c$，我们将向量 $u - cv$ 拆分为相互垂直的两部分：
              \begin{align*}
                  u - cv        & = (u - \text{proj}_{v}(u)) + (\text{proj}_{v}(u) - cv) \\
                  u_{\perp}     & = u - \text{proj}_{v}(u) \quad (\text{垂直于 } v)         \\
                  u_{\parallel} & = \text{proj}_{v}(u) - cv \quad (\text{平行于 } v)
              \end{align*}

              因为 $u_{\perp}$ 与 $u_{\parallel}$ 相互垂直，我们可以构建一个直角三角形。
              由勾股定理可知：
              $ \|u - cv\|^2 = \|u_{\perp}\|^2 + \|\text{proj}_{v}(u) - cv\|^2 $

              分析上式以求 $\|u - cv\|^2$ 的最小值：
              \begin{itemize}
                  \item 第一项 $\|u_{\perp}\|^2$ 为常数（不依赖于 $c$）；
                  \item 第二项 $\|\text{proj}_{v}(u) - cv\|^2 \ge 0$。
              \end{itemize}
              因此，当且仅当第二项为 $0$ 时，$\|u - cv\|^2$ 取得最小值。即：
              \begin{align*}
                  \text{proj}_{v}(u) - cv                                                & = 0 \\
                  \frac{\langle u, v \rangle}{\langle v, v \rangle} v - cv               & = 0 \\
                  \left( \frac{\langle u, v \rangle}{\langle v, v \rangle} - c \right) v & = 0
              \end{align*}
              设 $v \neq 0$，则有：
              $$ c = \frac{\langle u, v \rangle}{\langle v, v \rangle} $$
              此时，最小距离的平方为 $\|u - cv\|^2 = \|u_{\perp}\|^2$, $u - cv = u_{\perp} $。证毕。
    \end{itemize}
\end{flushleft}


\section{子空间投影}
\begin{flushleft}
    \begin{itemize}
        给定一个 k 维子空间 $W \subseteq V$, 把 $u$ 投影到 W 上. \\
        \textbf{问题}: 找到 $w \subseteq W$ 使得 $u- w \perp W$ (即 $u - w$ 与 $W$ 中所有向量正交).
        如果 $W$ 有一组正交基 ${w_{1}, w_{2}, ... ,w_{k}}$ (即 $w_{i} \perp w_{j}$ 当 $i \neq j$), 那么解非常干净:
        \begin{align}
            proj_w(u) = \sum_{i=1}^{k}\frac{<u, w_{i}>}{<w_{i},w_{i}>}w_{i}
        \end{align}

        为什么这是对的? 证明:

        设 $w$ 是存在的 \\
        那么它存在下列性质: \\
        $w \subseteq W \subseteq V$ w 是 V 里的向量，并且它落在子空间 W 中 \\
        w 是向量 那么我们可以代数式表达: \\
        \begin{align}
            w = c{1}w_{1} + c{2}w_{2} + ... c_{k}w_{k} \\
            w = \sum_{i=1}^{k}c_{i}w_{i}
        \end{align}

        投影要求: \\
        $u - w \perp W$ \\
        那么 $w_{j}$ 是 W 里任意一个基方向, 使用下列成立: \\
        \begin{align}
            <u - w, w_{j}> = 0                                \\
            <u - \sum_{i=1}^{k}c_{i}w_{i}, w_{j}> = 0         \\
            <u, w_{j}> - \sum_{i=1}^{k}c_{i}<w_{i},w_{j}> = 0 \\
        \end{align}
        因为 W 这些基都是正交的($i \ne j $)
        那么:
        \begin{align}
            <u, w_{j}> - \sum_{i=1}^{k}c_{i}<w_{i},w_{j}> = 0 \\
            <u, w_{j}> - c_{j}<w_{j},w_{j}> = 0               \\
            <u, w_{j}>  =   c_{j}<w_{j},w_{j}>                \\
            \frac{<u, w_{j}>}{<w_{j},w_{j}>}  = c_{j}
        \end{align}
        把 $c_{j}$ 换成标准下标 $c_{i}$  \\
        $ c_{i} = \frac{<u, w_{i}>}{<w_{i},w_{i}>}$ \\
        那么 w: \\
        $w = proje_{w}(u) = \sum_{i=1}^{k}\frac{<u, w_{i}>}{<w_{i},w_{i}>}w_{i}$ \\

        这里我们能明白一个很有意思的点 ${w_{j}}$ 因为是存在投影点 w  和 u 有正交性
        另一方面，由于基向量之间两两正交，\(\langle w_i,w_j\rangle\) 在 \(i\neq j\) 时为 0。
        因此当我们用 \(w_j\) 这个方向去检测投影点 \(p\) 时，只有 \(w_j\) 方向自己的分量会留下，
        其他方向的分量全部被正交性消掉。
        所以在正交基下，每个系数 \(c_j\) 可以被单独计算，不同方向之间没有耦合。

        如果 W 中的基不正交 那么 $<u,w_{j}> - \sum_{i=1}^{k}c_{i}<w_{i},w_{j}> = 0$ 就无法化简 需要确定计算出解
        那么最终的 关于 w 的 投影点公式是不成立的

    \end{itemize}
\end{flushleft}


\section{Gram-Schmidt: 正交化算法}
\begin{flushleft}
    \begin{itemize}

        输入一组线性无关的向量 ${v_{1}, v_{2}, ..., v_{k}}$ \\
        输出一组互相正交的向量 ${w_{1}, w_{2}, ..., w_{k}}$, 张成同一个子空间 \\
        重点: 输入组中的向量必须要互相线性无关 $v_{i} != cv_{j}$ 也就是每一个向量都有独立的方向 \\

        这里我们研究下 $v_{1},v_{2}$ \\
        我们已知 $v_{1}$ 和 $v_{2}$ 是线性无关 \\
        那么我们 可以对 $v_{1}, v_{2}$ 做正交分解和投影 \\
        \begin{align}
            w_{1} = v_{1}                                                                            \\
            v_{2} = \vec{w_{1}}_{\parallel} + \vec{w_{1}}_{\perp}                                    \\
            \vec{w_{1}}_{\parallel} = proj_{w_{1}}(v_{2}) = \frac{<v_{2},w_{1}>}{<w_{1},w_{1}>}w_{1} \\
            \vec{w_{1}}_{\perp} =  v_{2} -  \frac{<v_{2},w_{1}>}{<w_{1},w_{1}>}w_{1}                 \\
            w_{2} = \vec{w_{1}}_{\perp}                                                              \\
            <w_{1}, w_{2} > = 0                                                                      \\
        \end{align}

        那么这里我们得到了 $w_{1},w_{2}$ 它们是正交 \\
        那么我们思考下 $w_{3}$ 如果它和 $w_{2}$ 正交 那么它和 $w_{1}$ 的几何性质呢 ?  不能确定 \\
        但是我们需要保证 $w_{3} \perp w_{1}$  \\
        那么我可以得到一个概念 $w_{3} \perp w_{1} \&\& \; w_{3} \perp w_{2}$, $w_{3} \perp span(w_{1}, w_{2})$ \\
        从维度的角度来看 新延伸出来的的维度 新延伸出来的维度不一定天然正交旧维度 : Gram-Schmidt 是主动把它改造成正交旧子空间的新方向. \\
        这里我们要解决一个关键的点: 延续出来的方向是耦合的还解耦的  例如 : $w_{3}$ 是个方向是直角三角形的边 那么它拆解出来的两个向量(平行和垂直) 是否和 w1,w2耦合

        \textbf{1: 什么叫耦合什么叫解耦 ?}: \\
        我们设在3维空间中的子平面空间:
        \begin{align}
            W_{2} = span(w1,w2)       \\
            \vec{w1} = \begin{bmatrix}
                           1 \\
                           0 \\
                           0
                       \end{bmatrix} \\
            \vec{w2} = \begin{bmatrix}
                           0 \\
                           1 \\
                           0
                       \end{bmatrix} \\
        \end{align}
        $\vec{v3} = \begin{bmatrix}
                c1 \\
                c2 \\
                c3
            \end{bmatrix}$
        $ = \begin{bmatrix}
                c1 \\
                c2 \\
                0
            \end{bmatrix} (\text{as A}) + \begin{bmatrix}
                0 \\
                0 \\
                c3
            \end{bmatrix} (\text{as B})$ \\
        $A \in span(w1,w2)$, B 是延展方向 \\
        这里会发现一个有意思的点 $<A,B> = 0$ 也就是说 A, B 是正交的  \\
        我们把 B 转换成代数模式: $B = 0w1 + 0w2 + c3w3$ 代数式非常的干净的表现 B 的分量构成 w1,w2 占比为0 这是代数视角的解耦\\
        这里有一个非常重要的点 B 这个向量与 A 向量(或者更准确的说是 span(w1,w2)) 的几何形态 是完全垂直的 而非是斜的 \\
        B 无法分解出 $A \parallel + A\perp$ 或者说 $A \parallel = 0 $ 这是几何视角的解耦 \\
        对 A, B (向量) 来说, A 是耦合的, B 是解耦的(影响方向的分量只有一个) \\

        什么是耦合?
        $\vec{v3}$ 是一个耦合向量,
        起始点 $\vec{P}$, t 是常量 \\
        \begin{align}
            A \neq 0, B \neq 0      \\
            A \subseteq span(w1,w2) \\
            P_{new} = P + tv3       \\
            P_{new} = P + tA + tB
        \end{align}
        那么 我们很明显的就可以发现移动 "路径" 它必须在 $span(w1,w2)$ 先移动一段距离
        这也就是 v3 耦合了$span(w1,w2)$

        这里有一个很有意思的点
        我们找到 w3 使其正交于 $span(w_{1}, w_{2})$ 平面 在前面我们已经证明了当一个向量正交于一个子空间 那么它就正交与这个空间的所有基向量 \\


        \textbf{2: 剥离 v3 中的耦合 分解出 w3 使其正交 $span(w_{1}, w_{2})$}:
        我们已知 如果 v3 不正交 $span(w_{1},w_{2})$
        那么 $v3 = A + B, A \in span(w_{1},w_{2})$ 简单来说 也就是必须现在 $span(w_{1},w_{2})$ 移动一段距离
        那么 $w3 = v3 - A = B$
        那么存在下面代数式:
        \begin{align}
            W_{2} = span(w_{1}, w_{2}) \\
            A  = proj_{W_{2}}(v3)      \\
            w3 = v3 - proj_{W_{2}}(v3) \\
            w3 = v3 - proj_{w_{1}}(v3) - proj_{w_{2}}(v3); (w_{1} \perp w_{2})
        \end{align}
        这里有一个非常重要的点 w3 是因为正交 $W_{2}$ 子空间 并且 $w3 \notin W_{2}$, 才推导出 正交 $W_{2}$ 中所有的基向量

        \textbf{一般的: }
        \begin{align}
            w_{i} = v_{i} - \sum_{j=1}^{i-1}\frac{<v_{i},w_{j}>}{<w_{j},w_{j}>}w_{j}
        \end{align}

    \end{itemize}
\end{flushleft}

\end{document}